%% LyX 2.4.4 created this file.  For more info, see https://www.lyx.org/.
%% Do not edit unless you really know what you are doing.
\documentclass[english]{article}
\usepackage[T1]{fontenc}
\usepackage[latin9]{inputenc}
\usepackage{enumitem}
\usepackage{amsmath}
\usepackage{amsthm}
\usepackage{amssymb}
\usepackage{cancel}
\usepackage{geometry}
\geometry{verbose,tmargin=1in,bmargin=1in,lmargin=1in,rmargin=1in}
\PassOptionsToPackage{normalem}{ulem}
\usepackage{ulem}

\makeatletter

%%%%%%%%%%%%%%%%%%%%%%%%%%%%%% LyX specific LaTeX commands.
%% Because html converters don't know tabularnewline
\providecommand{\tabularnewline}{\\}

%%%%%%%%%%%%%%%%%%%%%%%%%%%%%% Textclass specific LaTeX commands.
\numberwithin{equation}{section}
\numberwithin{figure}{section}
\newlength{\lyxlabelwidth}      % auxiliary length 
\theoremstyle{plain}
\newtheorem{thm}{\protect\theoremname}[section]
\theoremstyle{definition}
\newtheorem{xca}[thm]{\protect\exercisename}

%%%%%%%%%%%%%%%%%%%%%%%%%%%%%% User specified LaTeX commands.
\usepackage{fancyhdr}
\usepackage{lastpage}
\pagestyle{fancy}
\usepackage{enumitem}
\usepackage{ifthen}
\everymath=\expandafter{\the\everymath\displaystyle}
%\DeclareMathSizes{10}{10}{10}{10}
\usepackage{multicol}

\usepackage{tikz}
\usetikzlibrary{patterns}
\usetikzlibrary{plotmarks}
\usepackage{pgfplots}
\pgfplotsset{compat=newest}

\usepackage{calc}
\usepackage[nomessages]{fp}

\usepackage{datetime}

% Added by lyx2lyx
\setlength{\parskip}{\smallskipamount}
\setlength{\parindent}{0pt}

\makeatother

\usepackage{babel}
\providecommand{\exercisename}{Exercise}
\providecommand{\theoremname}{Theorem}

\begin{document}
\renewcommand{\author}{Dr.\,James Ashe}

\newcommand{\classNum}{335}

\newcommand{\classSec}{A}

\newcommand{\nameType}{solutions}

\newcommand{\examType}{Homework 4}

%\newcommand{\Date}{\formatdate{6}{2}{2014}} %%\AdvanceDate[12] \today

%\newdate{my_date}{\day}{\month}{\year}%   
%\AdvanceDate[-5] 
%\displaydate{my_date}

%%First page's header%%%%%%%%%%%%%%%%%%%%%%
\fancypagestyle{firststyle} %%header settings for first pagr
{
	\fancyhead[L]{MTH \classNum{}(\classSec{}) \author{}\\
	\textbf{\examType{}}} %\textbf{\examType{}} \Date{}}
	\fancyhead[C]{\textbf{\nameType}}
	\setlength{\headsep}{14pt}
}
%%%%%%%%%%%%%%%%%%%%%%%%%%%%%%%%%%%%%%%%%%

%%Next page's header%%%%%%%%%%%%%%%%%%%%%%%
%\setlist{nolistsep} %eliminates space between list items
\lhead{MTH \classNum{}(\classSec{})} 
\chead{\textbf{\nameType}} 
\cfoot{\thepage{}~of~\pageref{LastPage}} 
\setlength{\headsep}{5pt} 
%%%%%%%%%%%%%%%%%%%%%%%%%%%%%%%%%%%%%%%%%%%

%%Various settings; see the preamble for other settings%%%%%
%%\setenumerate[0]{leftmargin=12pt,itemindent=0pt,labelwidth=0pt}
\renewcommand{\labelenumi}{\textbf{\arabic{enumi}.}}
\renewcommand{\labelenumii}{\textbf{\alph{enumii}.}}
\thispagestyle{firststyle}
%%%%%%%%%%%%%%%%%%%%%%%%%%%%%%%%%%%%%%%%%%

\null

%\vspace{1cm}

\setcounter{section}{3}

\Large{\textbf{Chapter 4}} \large

\textbf{\textbf{\examType} due: \formatdate{10}{4}{2014}}

All non-starred problems from chapter 4 section 1.

\setcounter{thm}{0}

\null

%\vspace{1cm}
\begin{xca}
\label{ex:Z is a group under +}Show that the set of integers, $\mathbb{Z}$,
is a group under addition. 

\begin{proof}
Any integer $n$ can be written as $n=1+1+\cdots1$ for $n$ $1$'s.
So then for any two $m,n\in\mathbb{Z}$ we have 
\begin{eqnarray*}
m+n & = & (1+1+\cdots+1)+(1+1+\cdots+1)\\
 & = & 1+1+\cdots+1
\end{eqnarray*}
which is just $m+n$ $1$'s. So $\mathbb{Z}$ is closed.

$1+0=0+1=1$ so then $n+0=1+1+\cdots+1+0=n$, and $0$ is the identity.

$1+(-1)=(-1)+1=0$. So then the inverse of $n$ is just $-n=-1-1\cdots-1$. 

Associativity follows since addition of any multiple of integers is
just adding a list of $1's$, and so the order is irrelevant. Thus
$(m+n)+m'=m+(n+m')$ for any $n,m,m'\in\mathbb{Z}$.
\end{proof}
\end{xca}

\begin{xca}
Prove that the set of even integers is a group under addition. 

\begin{proof}
Let $2\mathbb{Z}=\left\{ n\in\mathbb{Z}\mbox{ such that }2|n\right\} $
for $2$. Let $a,b\in2\mathbb{Z}$. Then $a=2x$ and $b=2y$ for some
$x,y\in\mathbb{Z}$. $a+b=2x+2y=2\left(x+y\right)$ so $2|a+b$, and
$2\mathbb{Z}$ is closed. The identity is $0$, $a^{-1}=-a=-2x$ which
is divisible by $2$, and associativity follows from associativity
of $\mathbb{Z}$.
\end{proof}
\end{xca}

\begin{xca}
Prove that integers of $n$ multiples for $n\in\mathbb{Z}^{+}$ are
a group under addition.

\begin{proof}
Let $n\mathbb{Z}=\left\{ x\in\mathbb{Z}\mbox{ such that }n|x\right\} $
for some integer $n$. Let $a,b\in n\mathbb{Z}$. Then $a=nx$ and
$b=ny$ for some $x,y\in\mathbb{Z}$. $a+b=nx+ny=n\left(x+y\right)$
so $n|a+b$, and $n\mathbb{Z}$ is closed. The identity is $0$, $a^{-1}=-a=-nx$
which is divisible by $n$, and associativity follows from associativity
of $\mathbb{Z}$.
\end{proof}
\end{xca}

\begin{xca}
\label{matrices are a group}Prove that the set of all $2\times2$
matrices with determinant $1$ and entries from $\mathbb{R}$ is a
group under matrix multiplication. 

\begin{proof}
We'll prove this for the more general case for $n\times n$ matrices.
Let $M_{i,j}$ represent the entry in the matrix $M$ at row $i$
and column $j$. Multiplication of matrices $A$ and $B$ is defined
by 
\[
(AB)_{i,j}={\displaystyle \sum_{k=1}^{m}A_{i,k}B_{k,j}}
\]
Where $m$ is the number of rows in $B$. This product is defined
if and only is the number of columns in $A$ equals the number of
rows in $B$. This will always be the case if both $A$ and $B$ are
$n\times n$ matrices.

If was previously show that $\mathbb{R}$ is a group under multiplication
and addition. So the each entry in $AB$, $AB_{i,j}$, is a real number.
Furthermore from the formula, we can see that if $A$ and $B$ are
$n\times n$ matrices then $AB$ is also a $2\times2$ matrix. Finally,
$A$ and $B$ are square matrices and $\det(A)=\det(B)=1$. This implies
that $\det(AB)=\det(A)\det(B)=1$, and thus this set is closed under
multiplication..

The identity will be the identity matrix, 
\[
I=\begin{bmatrix}1 & 0 & \cdots & 0\\
0 & 1 &  & 0\\
\vdots &  & \ddots & \vdots\\
0 & 0 & \cdots & 1
\end{bmatrix}
\]
As 
\begin{eqnarray*}
(AI)_{i,j} & = & {\displaystyle \sum_{k=1}^{m}A_{i,k}I_{k,j}=}A_{i,1}0+A_{i,2}0+\cdots+A_{i,j}1+\cdots+A_{n,m}0\\
(IA)_{ij} & = & {\displaystyle \sum_{k=1}^{m}I_{k,j}A_{i,k}=}0A_{i,1}+0A_{i,2}+\cdots+1A_{i,j}+\cdots+0A_{n,m}\\
 & = & A_{i,j}
\end{eqnarray*}

For each matrix $A$ there exists an $A^{-1}$ such that $AA^{-1}=A^{-1}A=I$
since $\det(A)=1$. 

For associativity 
\begin{eqnarray*}
(AB)_{i,j}C_{ij} & = & \left({\displaystyle \sum_{k=1}^{m}A_{i,k}B_{k,j}}\right)C_{ij}={\displaystyle \sum_{k=1}^{m}\left({\displaystyle \sum_{k=1}^{m}A_{i,k}B_{k,j}}\right)C_{k,j}}\\
 & = & {\displaystyle \sum_{k=1}^{m}{\displaystyle \sum_{k=1}^{m}}\left({\displaystyle A_{i,k}B_{k,j}}\right)C_{k,j}}\mbox{ by distribution}\\
 & = & {\displaystyle \sum_{k=1}^{m}{\displaystyle \sum_{k=1}^{m}}A_{i,k}\left(B_{k,j}C_{k,j}\right)}\mbox{ by associativity of \ensuremath{\mathbb{R}}.}\\
 & = & A_{i,j}(B_{i,j}C_{i,j})
\end{eqnarray*}
\end{proof}
\end{xca}

\begin{xca}
Show that the following are \emph{not} groups under the given operation.

\begin{enumerate}
\item $\mathbb{Z}$ under multiplication.

\begin{proof}
$2\in\mathbb{Z}$, but $2^{-1}\notin\mathbb{Z}$ since $2^{-1}2=1\Rightarrow2^{-1}=\frac{1}{2}$. 
\end{proof}
\item $S=\left\{ z|z\in\mathbb{Z}\mbox{ and }z=\sqrt{5}\right\} $ under
any operation.

\begin{proof}
$\sqrt{5}\notin\mathbb{Z}$ so then $S=\emptyset$. So $S$ is not
a group.
\end{proof}
\item $\mathbb{Q}$, the set of rationals under multiplication.

\begin{proof}
$0\in\mathbb{Q}$, but $0$ has no multiplicative inverse. 
\end{proof}
\item The set of odd integers under addition.

\begin{proof}
$3$ and $5$ are odd, but $3+5=8=2\cdot4$ which is not odd. So it
is not closed.
\end{proof}
\end{enumerate}
\end{xca}

\begin{xca}
Let $S=\left\{ a+b\sqrt{5}|\,a,b\mbox{ are not both zero and }a,b\in\mathbb{Q}\right\} $.
Prove that $S$ is a group. (Hint: rationalize the denominator to
show that the inverse of $a+b\sqrt{5}$ is in $S$.)

\begin{proof}
Let $x,y\in S$ so then $x=a+b\sqrt{5}$ and $y=c+d\sqrt{5}$ for
some $a,b,c,d\in\mathbb{Q}$. $xy=ac+\left(b+d\right)\sqrt{5}+5bd=ab+5bd+\left(a+b\right)\sqrt{5}$
which is in $S$ as $\mathbb{Q}$ is closed under multiplication and
addition. $1$ is the identity. $x^{-1}=\frac{1}{a+b\sqrt{5}}\frac{\left(a-b\sqrt{5}\right)}{\left(a-b\sqrt{5}\right)}=\frac{a-b\sqrt{5}}{a^{2}-5b}=\frac{a}{a^{2}-5b}+\frac{-b}{a^{2}-5b}\sqrt{5}\in S$.
Associativity follows by the associativity of $\mathbb{R}$. 
\end{proof}
\end{xca}

\begin{xca}
Show that $\left\{ 1,2,3\right\} $ under multiplication modulo $4$
is \emph{not }a group but that $\left\{ 1,2,3,4\right\} $ under multiplication
modulo $5$ is a group. (Example: $2\cdot3=6$ and $6=2\bmod4$ so
$[6]=[2]$ and thus $2\cdot3$ \emph{is} in $\left\{ 1,2,3\right\} $).

\begin{proof}
$2\cdot2=0\bmod4$ and $0\notin\left\{ 1,2,3\right\} $ so it is not
a group. 

$\left\{ 1,2,3,4\right\} $: %
\begin{tabular}{c|cccc}
 & 1 & 2 & 3 & 4\tabularnewline
\hline 
1 & 1 & 2 & 3 & 4\tabularnewline
2 & 2 & 4 & 1 & 3\tabularnewline
3 & 3 & 1 & 4 & 2\tabularnewline
4 & 4 & 1 & 2 & 3\tabularnewline
\end{tabular} 

So it is closed, the identity is $1$, each element has an inverse,
and associativity follows by the associativity of modular multiplication.
\end{proof}
\end{xca}

\begin{xca}
Addition in $\mathbb{Z}_{n}$ is defined as $\left[a\right]+\left[b\right]=\left[a+b\right]$.
Multiplication in $\mathbb{Z}_{n}$ is defined as $\left[a\right]\cdot\left[b\right]=\left[a\cdot b\right]$.
For example, $\mathbb{Z}_{4}=\left\{ \left[0\right],\left[1\right],\left[2\right],\left[3\right]\right\} $
for convenience we write $\mathbb{Z}_{4}=\left\{ 0,1,2,3\right\} $.
$2+3=5$ and $5=1\bmod4$ so $\left[5\right]=\left[1\right]$. Show
that $\mathbb{Z}_{2}\times\mathbb{Z}_{3}$ is a group under modular
addition. 

\begin{proof}
Closure: If $(a,x),(b,y)\in\mathbb{Z}_{2}\times\mathbb{Z}_{3}$ then
$(a,x)+(b,y)=(a+b,x+y)$. $\mathbb{Z}_{2}$ and $\mathbb{Z}_{3}$
are both groups so $a+b\in\mathbb{Z}_{2}$ and $x+y\in\mathbb{Z}_{3}$
thus $(a+b,x+y)\in\mathbb{Z}_{2}\times\mathbb{Z}_{3}$.

Identity: The identity in both $\mathbb{Z}_{2}$ and $\mathbb{Z}_{3}$
is $0$. Thus $(a,b)+(0,0)=(a+0,b+0)=(a,b)$. Similarly $(0,0)+(a,b)=(a,b)$.
So $(0,0)$ is the identity in $\mathbb{Z}_{2}\times\mathbb{Z}_{3}$. 

Inverse: The inverses of $a$ and $b$ is 
\end{proof}
\end{xca}

\begin{xca}
Prove that the set of all $n\times n$ matrices with real entries
is an abelian group under \uline{addition}.

\begin{proof}
Let $M_{i,j}$ represent the entry in the matrix $M$ at row $i$
and column $j$. Let $A$ and $B$ be $n\times n$ matrices then the
entry in the $i^{\mbox{th}}$ row and $j^{\mbox{th}}$ column is found
as follows: 
\[
(A+B)_{i,j}=A_{i,j}+B_{i,j}
\]
$A+B$ is then well defined as long as they have matching dimension.
Since the real numbers are a group under addition with identity $0$
and $a^{-1}=-a$ for any $a\in\mathbb{R}$, the addition of entries
is closed, has identity $0$, inverse $-A_{ij}$, and is associative.
Also the addition of entries is abelian since $\mathbb{R}$ is abelian. 

Since closure, associativity, and commutativity hold for each entry
it holds for matrices with real entries. The identity is just a matrix
will all zeros, 
\[
\mathbf{0}=\begin{bmatrix}0 & 0 & \cdots & 0\\
0 & 0 &  & 0\\
\vdots &  & \ddots & \vdots\\
0 & 0 & \cdots & 0
\end{bmatrix}
\]
 and the inverse for a matrix $A$ is just $-A$, e.g., 
\[
\mbox{if }A=\begin{bmatrix}a_{1,1} & a_{1,2} & \cdots & a_{1,n}\\
a_{2,1} & a_{2,2} &  & a_{2,n}\\
\vdots &  & \ddots & \vdots\\
a_{n,1} & a_{n,2} & \cdots & a_{n,n}
\end{bmatrix}\mbox{ then }A^{-1}=\begin{bmatrix}-a_{1,1} & -a_{1,2} & \cdots & -a_{1,n}\\
-a_{2,1} & -a_{2,2} &  & -a_{2,n}\\
\vdots &  & \ddots & \vdots\\
-a_{n,1} & -a_{n,2} & \cdots & -a_{n,n}
\end{bmatrix}
\]
\end{proof}
\end{xca}

\begin{xca}
Prove that the set of all $2\times2$ matrices with determinant $1$
and entries from $\mathbb{R}$ is a non-abelian group under matrix
multiplication (we already know that it is a group by exercise \ref{matrices are a group}. 
\end{xca}

\begin{proof}
Let 
\[
A=\begin{bmatrix}0 & 0\\
1 & 1
\end{bmatrix}\mbox{ and }B=\begin{bmatrix}0 & 1\\
0 & 1
\end{bmatrix}
\]
then 
\[
AB=\begin{bmatrix}0 & 0\\
0 & 2
\end{bmatrix}\mbox{ and }BA=\begin{bmatrix}1 & 1\\
1 & 1
\end{bmatrix}
\]
So $AB\neq BA$, and this group is not abelian. 
\end{proof}

\begin{xca}
$\bigstar$ 
\end{xca}

\begin{xca}
Prove that $G$ is an abelian group if and only is $\left(a\oplus b\right)^{-1}=a^{-1}\oplus b^{-1}$
for all $a,b\in G$. 

\begin{proof}
$(\Rightarrow)$ Suppose that $G$ is abelian. $\left(a\oplus b\right)^{-1}=b^{-1}\oplus a^{-1}$
since $b^{-1}\oplus a^{-1}\oplus(a\oplus b)=0_{G}$ and inverses are
unique. Since $G$ is assumed to be abelian 
\[
\left(a\oplus b\right)^{-1}=b^{-1}\oplus a^{-1}=a^{-1}\oplus b^{-1}\mbox{.}
\]

$(\Leftarrow)$ Suppose that $(a\oplus b)^{-1}=a^{-1}\oplus b^{-1}$.
Similarly to above $(a\oplus b)^{-1}=b^{-1}\oplus a^{-1}$ so then 

\begin{eqnarray*}
(a\oplus b)\oplus(a\oplus b)^{-1} & = & 0_{G}\\
a\oplus b\oplus(a^{-1}\oplus b^{-1}) & = & 0_{G}\\
a\oplus b\oplus(a^{-1}\oplus b^{-1})\oplus(b\oplus a) & = & 0\oplus(b\oplus a)\\
a\oplus b\oplus(a^{-1}\oplus\cancel{(b^{-1}\oplus b)}\oplus a) & = & b\oplus a\\
a\oplus b\oplus\cancel{(a^{-1}\oplus a)} & = & b\oplus a\\
a\oplus b & = & b\oplus a
\end{eqnarray*}

This is easier to see using juxtaposition to denote the operation
instead of ``$\oplus$''. 

$(\Rightarrow)$ Suppose that $G$ is abelian. Then $(ab)^{-1}=(ba)^{-1}$.
$(ba)^{-1}ba=e\Rightarrow(ba)^{-1}=a^{-1}b^{-1}$ as required.

$(\Leftarrow)$ Suppose that $(ab)^{-1}=a^{-1}b^{-1}$. Similarly
to above $(ab)^{-1}=b^{-1}a^{-1}$ So then $ab(ab)^{-1}=e\Rightarrow ab(a^{-1}b^{-1})=e\Rightarrow abab^{-1}ba=ba\Rightarrow ab=ba$ 
\end{proof}
\end{xca}

\begin{xca}
Suppose that if $G$ is a group such that if $a,b,c\in G$ and $a\oplus b=c\oplus a$
then $b=c$. Prove that $G$ is an Abelian group, i.e., that $a\oplus b=b\oplus a$
for all $a,b\in G$. Hint: consider the element $a\oplus(b\oplus a)=(a\oplus b)\oplus a$.

\begin{proof}
Let $a,b,c\in G$ and suppose that $a\oplus b=c\oplus a$ implies
$b=c$. $G$ is a group so then $a\oplus(b\oplus a)=(a\oplus b)\oplus a$
by associativity. By the assumption $a\oplus(b\oplus a)=(a\oplus b)\oplus a$
implies that $a\oplus b=b\oplus a$ so then $G$ is abelian.
\end{proof}
\end{xca}


\end{document}
